\section{本章小结}
\label{sec:dapx-summary}

章首的问题是：给定一个只能逐点计算未归一化密度的分布，怎样找一个容易计算的分布代替它。本章把这个问题拆成四个选择，每个选择都有自己的依据和代价。表~\ref{tab:dapx-summary}把它们并列。

\begin{BookInlineTable}
 \centering
 \caption{逼近问题中的四个选择。}
 \label{tab:dapx-summary}
 \begin{tabular}{lll}
  \toprule
  选择 & 选项 & 依据与代价\\
  \midrule
  候选集合 & 乘积结构 & 删去组间依赖，因子形状不限\\
   & 有限维参数族 & 限制形状，参数须保持合法\\
  KL方向 & $D_{\mathrm{KL}}(Q\Vert\Pi)$ & 需$Q\ll\Pi$；只需$\widetilde{\pi}$与$Q$下期望\\
   & $D_{\mathrm{KL}}(\Pi\Vert Q)$ & 需$\Pi\ll Q$；需$\Pi$下期望或样本\\
  求解 & 逐因子更新 & 子问题有闭式解，目标值单调下降\\
   & 参数梯度 & 适用于参数族，需梯度与步长\\
  评价 & 误差分解 & 表达限制、优化差距、计算误差\\
   & 按任务检验 & 矩、事件概率；无界函数要尾部条件\\
  \bottomrule
 \end{tabular}
\end{BookInlineTable}

候选集合决定了能保留什么。乘积候选无论怎样选择，都不能表示组间的依赖：概率表上两个变量常常相等的事实、相关高斯中的协方差都会消失。有限维参数族限制的是形状，单个高斯不能有两个峰，也不能像四次方目标那样有比高斯更轻的尾。这些损失在逼近开始之前就已确定，即式~\eqref{eq:dapx-error-parts}中的表达限制。

KL方向决定了在这些损失之间怎样取舍。按候选加权的方向惩罚候选在目标稀少处放置的质量，只要求逐点计算$\widetilde{\pi}$，因此在目标难算时最常用；它的最优解在三状态例子里给中间状态多分了概率，在相关高斯中把方差压到条件方差$1-\rho^{2}$，在双峰目标上贴住一个峰。按目标加权的方向惩罚候选对目标的覆盖不足，对指数族候选就是矩匹配，却需要目标下的期望。两个方向选出的分布一般不同，“择峰”与“覆盖”只是倾向，具体结论要按例子计算。

求解方法由候选集合的类型决定。乘积候选且因子不受限制时，每个因子的最优更新是对其余因子平均目标的对数密度，再取指数并归一化；每次更新不增加目标值，但只保证目标值收敛，可能停在鞍点或局部最优。参数族候选则用梯度下降，重新参数化保证每一步仍是合法分布，期望没有闭式时需要随机梯度。

评价时要分清误差的来源，并直接检验任务关心的量。四次方目标上，表达限制是$0.0472$，$10$步梯度下降后的优化差距是$0.018$，而用$100$个样本估计目标值的误差约为$0.043$，足以掩盖优化差距。两个方向的最优高斯分别保留了四阶矩和二阶矩，各自在另一个矩上有明显误差。扩大候选集合只降低表达限制的下确界，实际结果还取决于优化能否找到好的解。

本章假定目标分布已经给定。在统计推断中，后验分布常常正是这样一个只知道未归一化密度的目标，本章的方法可以直接用于计算它的近似；后验分布怎样由先验与数据建立、近似结果在统计上意味着什么，见\BookChapterRef{chap:parametric-statistical-inference}{卷八第1章“参数化统计推断”}。参数空间中同样长的一步引起的分布变化不同，怎样按分布的变化来度量并选择更新方向，由第~\ref{chap:information-geometry}章研究。不用近似分布、直接从目标抽样来计算期望的方法，以及这些期望估计的误差与方差，由第~\ref{chap:sample-generation}章至第~\ref{chap:expectation-gradient-estimation}章研究。
